\begin{tcolorbox}[
    colback=SkyBlue!5,
    colframe=blue!20,
    boxrule=1pt,
    arc=2pt,
    width=\textwidth,
    breakable
]

\begin{example}
\cref{fig:r+w+w-knuth} shows $\mathcal{D}(\ref{ex:r+w+w})$. 
Notice that the tree is \emph{different} from $\mathcal{T}(P)$,
even though both are trees.
When we run \textbf{Algorithm T}, it picks the leaves 
with probabilities $1/6$, $1/6$, $1/12$, $1/12$, $1/4$, and $1/4$,
respectively.
The expectation is 6, which is equal to $C(P)$.

For this example, \textbf{Algorithm T} has higher variance than running 
\textbf{Algorithm K} on $\mathcal{T}(P)$.
However, the advantage is that \textbf{Algorithm T} works on programs for which
$\mathcal{T}(P)$ is not a tree.
For example, \textbf{Algorithm T} on both the programs \ref{ex:r+r+r} and \ref{ex:r+rr} returns 1. For both of these programs, their respective trees, $\mathcal{D}$(\ref{ex:r+r+r}) and $\mathcal{D}$(\ref{ex:r+rr}) contain exactly one
path and so the variance is 0. 
\qed
\end{example}
\centering
\begin{tikzpicture}[
    >=stealth,
    node distance=1.4cm and 1.8cm,
    event/.style={draw=none, font=\normalsize},
    rf/.style={->, thick, dashed, green!60!black},
    co/.style={->, thick, dotted, orange},
    hyper edge/.style={->, very thick, blue!70!black},
    graph box/.style={draw, rectangle, rounded corners=4pt, very thick, minimum width=1.4cm, minimum height=1cm, inner sep=0.1cm, minimum size=0cm},
graph box yellow/.style={graph box, dashed},
graph box blue/.style={graph box},
    scale=0.31,
    every node/.style={transform shape},
    edge-t1/.style={->, very thick, lime!60!black},
    edge-t2/.style={->, very thick, red!70!black},
    circ-label-t1/.style={text=black, inner sep=1pt, font=\normalsize, scale=2.5},
circ-label-t2/.style={text=black, inner sep=1pt, font=\normalsize, scale=2.5}
]

\pgfdeclarelayer{background}
\pgfsetlayers{background,main}

%% Level 0
% G1
\begin{scope}[xshift=0, yshift=0, local bounding box=G1box]
  \node[event] (g1i) {init};
  \node[event, below=0.7cm of g1i, xshift=-1.5cm] (g1r1) {$\opr(x)$};
  \draw[->] (g1i) -- node[midway, right, font=\small] {} (g1r1);
  \draw[rf] (g1i) to[bend right=20] (g1r1);
\end{scope}
\begin{pgfonlayer}{background}
\node[graph box yellow, fit=(G1box)] (G1) {};
\end{pgfonlayer}

%% Level 1
% G2
\begin{scope}[xshift=-8.5cm, yshift=-3cm, local bounding box=G2box]
  \node[event] (g2i) {init};
  \node[event, below=0.7cm of g2i, xshift=-1.5cm] (g2r1) {$\opr(x)$};
  \node[event, below=0.7cm of g2i] (g2w2) {$\opw(x)$};
  \draw[->] (g2i) -- node[midway, right, font=\small] {}(g2r1);
  \draw[->] (g2i) -- node[midway, right, font=\small] {}(g2w2);
  \draw[rf] (g2i) to[bend right=20] (g2r1);
\end{scope}
\begin{pgfonlayer}{background}
\node[graph box yellow, fit=(G2box)] (G2) {};
\end{pgfonlayer}

% G8
\begin{scope}[xshift=10cm, yshift=-3cm, local bounding box=G8box]
  \node[event] (g8i) {init};
  \node[event, below=0.7cm of g8i, xshift=-1.5cm] (g8r1) {$\opr(x)$};
  \node[event, below=0.7cm of g8i] (g8w2) {$\opw(x)$};
  \draw[->] (g8i) -- node[midway, left, font=\small] {}(g8r1);
  \draw[->] (g8i) -- node[midway, right, font=\small] {}(g8w2);
  \draw[rf] (g8w2) to[] (g8r1);
\end{scope}
\begin{pgfonlayer}{background}
\node[graph box yellow, fit=(G8box)] (G8) {};
\end{pgfonlayer}

%% Level 2
% G3
\begin{scope}[xshift=-14.3cm, yshift=-7cm, local bounding box=G3box]
  \node[event] (g3i) {init};
  \node[event, below=0.7cm of g3i, xshift=-1.5cm] (g3r1) {$\opr(x)$};
  \node[event, below=0.7cm of g3i] (g3w2) {$\opw(x)$};
  \node[event, below=0.7cm of g3i, xshift=1.5cm] (g3w3) {$\opw(x)$};
  \draw[->] (g3i) -- node[midway, left, font=\small] {}(g3r1);
  \draw[->] (g3i) -- node[midway, right, font=\small] {}(g3w2);
  \draw[->] (g3i) -- node[midway, right, xshift=0.3cm, font=\small] {}(g3w3);
  \draw[rf] (g3i) to[bend right=20] (g3r1);
  \draw[co] (g3w2) -- (g3w3);
\end{scope}
\begin{pgfonlayer}{background}
\node[graph box blue, fit=(G3box)] (G3) {};
\end{pgfonlayer}

% G4
\begin{scope}[xshift=-9.2cm, yshift=-7cm, local bounding box=G4box]
  \node[event] (g4i) {init};
  \node[event, below=0.7cm of g4i, xshift=-1.5cm] (g4r1) {$\opr(x)$};
  \node[event, below=0.7cm of g4i] (g4w2) {$\opw(x)$};
  \node[event, below=0.7cm of g4i, xshift=1.5cm] (g4w3) {$\opw(x)$};
  \draw[->] (g4i) -- node[midway, left, font=\small] {}(g4r1);
  \draw[->] (g4i) -- node[midway, right, font=\small] {}(g4w2);
  \draw[->] (g4i) -- node[midway, right, xshift=0.3cm, font=\small] {}(g4w3);
  \draw[rf] (g4i) to[bend right=20] (g4r1);
  \draw[co] (g4w3) -- (g4w2);
\end{scope}
\begin{pgfonlayer}{background}
\node[graph box blue, fit=(G4box)] (G4) {};
\end{pgfonlayer}

% G5
\begin{scope}[xshift=-4cm, yshift=-7cm, local bounding box=G5box]
  \node[event] (g5i) {init};
  \node[event, below=0.7cm of g5i, xshift=-1.5cm] (g5r1) {$\opr(x)$};
  \node[event, below=0.7cm of g5i, xshift=1.5cm] (g5w3) {$\opw(x)$};
  \draw[->] (g5i) -- node[midway, left, font=\small] {}(g5r1);
  \draw[->] (g5i) -- node[midway, right, xshift=0.3cm, font=\small] {}(g5w3);
  \draw[rf] (g5w3) to[] (g5r1);
\end{scope}
\begin{pgfonlayer}{background}
\node[graph box yellow, fit=(G5box)] (G5) {};
\end{pgfonlayer}

%%% Level 3
% G6
\begin{scope}[xshift=-7cm, yshift=-11cm, local bounding box=G6box]
  \node[event] (g6i) {init};
  \node[event, below=0.7cm of g6i, xshift=-1.5cm] (g6r1) {$\opr(x)$};
  \node[event, below=0.7cm of g6i] (g6w2) {$\opw(x)$};
  \node[event, below=0.7cm of g6i, xshift=1.5cm] (g6w3) {$\opw(x)$};
  \draw[->] (g6i) -- node[midway, left, font=\small] {}(g6r1);
  \draw[->] (g6i) -- node[midway, right, font=\small] {}(g6w2);
  \draw[->] (g6i) -- node[midway, right, xshift=0.3cm, font=\small] {}(g6w3);
  \draw[rf] (g6w3) to[bend left=20] (g6r1);
  \draw[co] (g6w2) -- (g6w3);
\end{scope}
\begin{pgfonlayer}{background}
\node[graph box blue, fit=(G6box)] (G6) {};
\end{pgfonlayer}

% G7
\begin{scope}[xshift=-1cm, yshift=-11cm, local bounding box=G7box]
  \node[event] (g7i) {init};
  \node[event, below=0.7cm of g7i, xshift=-1.5cm] (g7r1) {$\opr(x)$};
  \node[event, below=0.7cm of g7i] (g7w2) {$\opw(x)$};
  \node[event, below=0.7cm of g7i, xshift=1.5cm] (g7w3) {$\opw(x)$};
  \draw[->] (g7i) -- node[midway, left, font=\small] {}(g7r1);
  \draw[->] (g7i) -- node[midway, right, font=\small] {}(g7w2);
  \draw[->] (g7i) -- node[midway, right, xshift=0.3cm, font=\small] {}(g7w3);
  \draw[rf] (g7w3) to[bend left=20] (g7r1);
  \draw[co] (g7w3) -- (g7w2);
\end{scope}
\begin{pgfonlayer}{background}
\node[graph box blue, fit=(G7box)] (G7) {};
\end{pgfonlayer}

% G9
\begin{scope}[xshift=5cm, yshift=-7cm, local bounding box=G9box]
  \node[event] (g9i) {init};
  \node[event, below=0.7cm of g9i, xshift=-1.5cm] (g9r1) {$\opr(x)$};
  \node[event, below=0.7cm of g9i] (g9w2) {$\opw(x)$};
  \node[event, below=0.7cm of g9i, xshift=1.5cm] (g9w3) {$\opw(x)$};
  \draw[->] (g9i) -- node[midway, left, font=\small] {}(g9r1);
  \draw[->] (g9i) -- node[midway, right, font=\small] {}(g9w2);
  \draw[->] (g9i) -- node[midway, right, xshift=0.3cm, font=\small] {}(g9w3);
  \draw[rf] (g9w2) to[] (g9r1);
  \draw[co] (g9w2) -- (g9w3);
\end{scope}
\begin{pgfonlayer}{background}
\node[graph box blue, fit=(G9box)] (G9) {};
\end{pgfonlayer}

% G10
\begin{scope}[xshift=13.5cm, yshift=-7cm, local bounding box=G10box]
  \node[event] (g10i) {init};
  \node[event, below=0.7cm of g10i, xshift=-1.5cm] (g10r1) {$\opr(x)$};
  \node[event, below=0.7cm of g10i] (g10w2) {$\opw(x)$};
  \node[event, below=0.7cm of g10i, xshift=1.5cm] (g10w3) {$\opw(x)$};
  \draw[->] (g10i) -- node[midway, left, font=\small] {}(g10r1);
  \draw[->] (g10i) -- node[midway, right, font=\small] {}(g10w2);
  \draw[->] (g10i) -- node[midway, right, xshift=0.3cm, font=\small] {}(g10w3);
  \draw[rf] (g10w2) to[] (g10r1);
  \draw[co] (g10w3) -- (g10w2);
\end{scope}
\begin{pgfonlayer}{background}
\node[graph box blue, fit=(G10box)] (G10) {};
\end{pgfonlayer}

%%%% Hyper Edges

\draw[edge-t1] (G1) -- node[above, yshift=8pt, circ-label-t1] {$\fwdtikz$} (G2);

\draw[edge-t2] (G1) -- node[above, yshift=8pt, circ-label-t2] {$\bwdtikz$} (G8);

\draw[edge-t1] (G2) -- node[left, yshift=20pt, circ-label-t1] {$\fwdtikz$} (G3);

\draw[edge-t1] (G2) -- node[right, xshift=8pt, circ-label-t1] {$\fwdtikz$} (G4);

\draw[edge-t2] (G2) -- node[right, yshift=20pt, circ-label-t2] {$\bwdtikz$} (G5);

\draw[edge-t1] (G8) -- node[left, yshift=20pt, circ-label-t1] {$\fwdtikz$} (G9);

\draw[edge-t1] (G8) -- node[right, yshift=20pt, circ-label-t1] {$\fwdtikz$} (G10);

\draw[edge-t1] (G5) -- node[left, yshift=5pt, xshift=-10pt, circ-label-t1] {$\fwdtikz$} (G6);

\draw[edge-t1] (G5) -- node[right, yshift=5pt, xshift=10pt, circ-label-t1] {$\fwdtikz$} (G7);

\end{tikzpicture}
\captionof{figure}{The transition system $\mathcal{D}(\ref{ex:r+w+w})$}
\label{fig:r+w+w-knuth}
\end{tcolorbox}